\documentclass[10pt,a4paper]{amsart}
\usepackage[margin=19mm]{geometry}
\usepackage{amsmath,amssymb,mathtools}
\usepackage[scr=rsfs,cal=euler]{mathalfa}
\usepackage{xfrac}
\usepackage[unicode,hidelinks,bookmarks=false]{hyperref}
\newcommand{\ie}{i.e.\ }
\newcommand{\cf}{cf.\ }
\newcommand{\resp}{respectively\ }
\newcommand{\Subsection}{\subsection}
\providecommand{\nobreakdash}{\nobreak\hbox{-}}
\setlength{\emergencystretch}{2em}
\multlinegap0pt
\usepackage{mathtools}
\usepackage{amssymb}
\usepackage{dsfont}
\usepackage{tikz}
\newtheorem{thm}{Theorem}
\newtheorem{lem}{Lemma}
\newtheorem{cor}{Corollary}
\def\A{\mathcal A} 
\def\B{\mathcal B} 
\newcommand{\End}{\operatorname{End}}
\newcommand{\Ext}{\operatorname{Ext}}
\newcommand{\Fun}{\operatorname{Fun}}
\def\Ga{\Gamma}
\newcommand{\Hom}{\operatorname{Hom}}
\def\Om{\Omega}
\def\bbK{\mathbb K} 
\newcommand{\fl}{\operatorname{fl}}
\newcommand{\frf}{\mathfrak{f}} 
\newcommand{\frt}{\mathfrak{t}}
\newcommand{\iso}{\xrightarrow{\raisebox{-.4ex}[0ex][0ex]{$\scriptstyle{\sim}$}}}
\newcommand{\lto}{\longrightarrow}
\renewcommand{\mod}{\operatorname{mod}}
\newcommand{\op}{\mathrm{op}}
\newcommand{\proj}{\operatorname{proj}} 
\newcommand{\rep}{\operatorname{rep}}
\newcommand{\xto}{\xrightarrow}
\title{Ordnung muss sein}
\author{Henning Krause}
\date{Source: arXiv:2602.01316v3 (CC BY 4.0)}
\begin{document}
\maketitle

\noindent\emph{Source and licence.} Ordnung muss sein. Version arXiv:2602.01316v3 (CC BY 4.0), distributed by arXiv under the Creative Commons Attribution 4.0 International licence, http://creativecommons.org/licenses/by/4.0/, as recorded in the arXiv licence metadata for this exact version and shown on the abstract page https://arxiv.org/abs/2602.01316v3. Full licence notice: \texttt{licenses/arxiv-2602.01316-CC-BY-4.0.txt}.\par\smallskip

\noindent\textit{Editorial scope.} Counted unit: the article's thm th:main2 (source0.tex lines 1066--1080). The article's complete original proof of it is delivered with it (source0.tex lines 1082--1136, one \texttt{proof} environment of 55 lines). No mathematics is written, repaired or re-derived: the statement and the proof are the article's own lines, in the article's own notation, and no retained line is substituted or re-flowed.  Retained uncounted as the source-local context the proof needs: lines 365--378; lines 570--574; lines 591--600; lines 624--632; lines 949--954; lines 985--991; lines 1033--1041.  Nothing was omitted from any retained span, so the package carries no omission record.  No retained line contains a citation, so the wrapper carries no bibliography and no bibliography entry.  No retained line contains a figure environment, an image-inclusion command or drawing code, so no image is delivered and the package declares no asset dependency.  Editorial formatting only, in addition: an AMS \texttt{amsart} preamble that restates the article's own declarations and macros used by the retained lines, namely the environments \texttt{thm}, \texttt{lem}, \texttt{cor} on separate counters, together with the standard packages those lines need.  The retained environments print in this document as Theorem 1, Lemma 1, Corollary 1 in order of appearance, whereas the article prints the same items with the numbering of its own full document.  The complete native source of the article as distributed in the arXiv e-print for this exact version is bundled unchanged as \texttt{sources/193782/source0.tex}.
\begin{thm}\label{th:main}
  A connected length category is equivalent to the category of
  pointwise finite dimensional representations of a poset over a field
  if and only if the following hold:
  \begin{enumerate}
  \item[(L1)] The ext-quiver  has no oriented cycle.
  \item[(L2)] There exists an object $M$ with a distributive lattice of
    subobjects such that each object is a subquotient of a finite
    direct sum of copies of $M$.
  \item[(L3)] For each simple object all its endomorphisms are central. 
\end{enumerate}
In this case the field equals the centre of the category and the poset
is given by the isomorphism classes of simple objects.
\end{thm}

\begin{lem}\label{le:max}
  Suppose the ext-quiver of $\A$ is acyclic. Let $M\to S$ be an essential
  epimorphism such that the canonical map $\End(M)\to\Ga$ is surjective,
 $\langle M\rangle=\A$, and $[M:S]=1$. Then $M$ is projective.
\end{lem}

\begin{lem}\label{le:torsion}
  Let $\Omega=\Omega'\sqcup \Omega''$ be a decomposition such that
  $\Ext^1(S_{x'}, S_{x''})=0$ for all $x'\in \Omega',x''\in
  \Omega''$. Then $(\A_{\Omega'},\A_{\Omega''})$ is a torsion pair for
  $\A$ and each object $X\in\A$ fits into an exact sequence
  \begin{equation*}\label{eq:tor}
    0\lto \frt_{\Omega'}(X)\lto X\lto \frf_{\Omega''}(X)\lto 0.
\end{equation*}    
Moreover, the functors $\frt_{\Omega'}$ and  $\frf_{\Omega''}$ are exact.
\end{lem}

\begin{lem}\label{le:centre}
  Let $\A$ be a length category such that the ext-quiver of $\A$ is
  acyclic and connected. Then the following are equivalent.
  \begin{enumerate}
  \item For each simple object all its endomorphisms are central.
  \item The centre of $\A$ is a field and isomorphic to the
    endomorphism ring of each simple object.
  \end{enumerate}
\end{lem}

\begin{cor}\label{co:ptfl}
A pointwise length category is determined by its
full subcategory of finite length objects. In particular, a functor
$\A\to\B$ between  pointwise length categories is an equivalence if it restricts
to an equivalence $\fl\A\iso\fl\B$ and preserves directed colimits.\qed
\end{cor}

\begin{lem}\label{le:ptfl}
  The assignment $M\mapsto\bar M$ induces a 
  fully faithful functor
  \[h\colon\A\lto\Fun_\bbK((\proj\A)^\op,\mod\bbK)\] that preserves
  colimits. Moreover, it identifies $\fl\A$ with the subcategory of
  finite length objects of $\Fun_\bbK((\proj\A)^\op,\mod\bbK)$.
\end{lem}

\begin{lem}\label{le:downfinite}
For $P$  and $\bbK$ the following are equivalent.
\begin{enumerate}
\item The poset $P$ is down-finite, that is,
  $\downarrow\! x$ is a finite set for each $x\in P$.
\item The category $\rep(P,\bbK)$ is a pointwise length category.
\item The representation $M_P$  is a sum of finite length subobjects. 
\end{enumerate}
\end{lem}

\begin{thm}\label{th:main2}
  A connected pointwise length category is equivalent to the category
  of pointwise finite dimensional representations of a poset over a
  field if and only if the following hold:
  \begin{enumerate}
  \item[(L1)] The ext-quiver admits only finitely many paths ending in a fixed vertex. 
  \item[(L2)] There is an object $M$ with a distributive lattice of
    subobjects such that each  finite length object is a subquotient of a finite
    direct sum of copies of $M$.
  \item[(L3)] For each simple object all its endomorphisms are central. 
\end{enumerate}
In this case the field equals the centre of the category and the poset is
given by the isomorphism classes of simple objects, with $[S]\le [T]$
if there is a path from $[S]$ to $[T]$ in the ext-quiver.
\end{thm}

\begin{proof}
  Fix a connected pointwise length category $\A$.  Suppose first that
  $\A\iso\rep(P,\bbK)$ for a poset $P$ and a field $\bbK$. Then $P$ is
  down-finite by Lemma~\ref{le:downfinite}. We need to show that the
  conditions (L1)--(L3) hold. The proof of Theorem~\ref{th:main}
  carries over. For instance, (L1) holds because the ext-quiver
  identifies with the Hasse diagram of $P$, and (L3) is clear as well.
  For (L2) we can reduce to the case that $P$ is finite, because any
  finite subset of $P$ is contained in a finite downward closed subset
  $P'\subseteq P$. Then $\A_{P'}$ is a length category and
  $\A_{P'}\iso\rep(P',\bbK)$. Moreover one uses that
  $\frt_{P'}(M_P)=M_{P'}$. Thus the object
  $M_P=\sum_{P'\subseteq P} M_{P'}$ is distributive since all
  finite length subobjects are distributive and finite intersections
  distribute over directed unions. It follows that (L2) holds for $\A$,
  given that
  \[\fl\A=\bigcup_{P'\subseteq P}\A_{P'}.\]

  Now suppose that conditions (L1)--(L3) hold for $\A$. First observe
  that \[\bbK\coloneqq Z(\fl\A)\cong Z(\A)\] is a field by
  Lemma~\ref{le:centre}. Choose a representative set of simple
  objects $(S_x)_{x\in \Om}$ in $\A$. Then $\Om$ is partially ordered  via $x\le y$ if
  there is a path from $S_x$ to $S_y$ in the ext-quiver of $\A$. 

  For $x\in \Om$ consider the subobject
  \[M_x\coloneqq \frt_{\downarrow x}(M)\subseteq M.\] We claim that $M_x$
  has finite length and that it is a projective cover of $S_x$. We
  apply Lemma~\ref{le:torsion} using the decomposition
  $\Om=\downarrow\!  x\sqcup(\Om\smallsetminus\downarrow\! x)$. Note that
  $\downarrow\!  x$ is a finite set because of (L1).
Thus $\A_{\downarrow x}$ is a length category and we wish to apply
  Theorem~\ref{th:main}.  It is easily checked that (L1)--(L3) hold; in
  particular $M_x$ is distributive and all objects of
  $\A_{\downarrow x}$ are subquotients of finite direct sums of $M_x$
  since the functor $\frt_{\downarrow x}$ is exact. Thus $M_x$ is
  projective in $\A_{\downarrow x}$ by Lemma~\ref{le:max},
  and then also projective in $\A$ since the inclusion
  $\A_{\downarrow x}\to\A$ is left adjoint to an exact functor.

  We have $\Hom(M_x,M_y)\cong\bbK$ if and only if $x\le y$, and a
  basis element is given by the inclusion $M_x\to M_y$. This follows
  as in the proof of Theorem~\ref{th:main} and yields a $\bbK$-linear
  equivalence $\bbK \Om\iso\proj\A$, where $\bbK \Om$ denotes the
  $\bbK$-linearisation of $\Om$.  Now consider the composite
\[\A\xto{\ h\ }\Fun_\bbK((\proj\A)^\op,\mod\bbK)\iso\Fun_\bbK((\bbK
  \Om)^\op,\mod\bbK) \iso\rep(\Om,\bbK)\]
where the first  functor is
from Lemma~\ref{le:ptfl}. This is fully faithful and we need to
show that it is essentially surjective. The functor sends $M$ to $M_\Om$
and this is a sum of finite length objects. Thus $\rep(\Om,\bbK)$ is a
pointwise length category by Lemma~\ref{le:downfinite}. The functor
identifies $\fl\A$ with the full subcategory of finite length objects
of $\rep(\Om,\bbK)$ by Lemma~\ref{le:ptfl}. Thus it follows from
Corollary~\ref{co:ptfl} that $\A\to \rep(\Om,\bbK)$ is an equivalence.
\end{proof}

\end{document}
